\scanpage{10}
et dans ces derniers homomorphismes, il n'est plus nécessaire de faire
mention de l'entier $n_{\unsure{o|0}}$ (i.e. ils ne dépendent pas du choix de l'entier
$n_o$ tel que $M\handwritten{\in}\mathrm{Ob}\,\underline C_{n_o}$\handwritten{)}. Pour notre objet, nous aurons plutot à considérer la situation duale d'un foncteur contravariant sur $\underline C$ :
\[
T:\underline C^o\handwritten{\longrightarrow}\underline C'\quad,
\]
donnant lieu à des homomorphismes
\[
T_{\deleted{n}}(M)\handwritten{\longrightarrow}\operatorname{Hom}_A(M,T_{\deleted{\unsure{n}}}(A_n))\qquad(n\handwritten{\geq}n_o)\quad,\tag*{(\%)}
\]
satisfaisant les propriétés de transitivité qu'il faut pour pouvoir passer
à la lim\deleted{o}ite (sous réserve d'existence des $\underset{\handwritten{\longrightarrow}}{\lim}$ en question)
\[
T(M)\handwritten{\unsure{-|\to}}\underset{\handwritten{\longrightarrow}}{\lim}_{\handwritten{n}}\mathbb H\mathrm{om}_A(M,T(A_n))\handwritten{\to}\operatorname{Hom}_A(M,\underset{\handwritten{\longrightarrow}}{\lim}_{\handwritten{n}}T(A_n))
\]
Pour que $T$ soit exact à gauche, il f et s que les $T_n$ le soient, ou encore
que les homomorphismes (\%) soient des isomorphismes. Si dans $\underline C'$ les
limites inductives \inserted{filtrantes} dénombrables existent et sont ``exactes'', à fortiori
si l'axiome (Ab 5) est vérifié, alors \deleted{l'exactitude à gauche de $T$ impliqque donc que l'homomorphisme} on voit fa\handwritten{cil}ement que le foncteur
$P\handwritten{\leadsto}\operatorname{Hom}_A(M,P)$ de \deleted{$\underline C'$ dans $\underline C'$} $A(\underline C')$ dans $\underline C'$ commute aux limites inductives \inserted{filtr} dénombrables (pour tout $M$ de présentation finie sur $A$), Donc \deleted{sous} si
\inserted{\handwritten{de plus}} $T$ est exact à gauche, l'homomorphisme composé
\[
T(M)\handwritten{\longrightarrow}\operatorname{Hom}_A(M,\underset{\handwritten{\longrightarrow}}{\lim}_{\handwritten{n}}T(A_n))
\]
est un isomorphisme, la réciproque étant triviale. Il y a lieu de
résumer ce résultat dans une

\uline{Proposition} 2.9.

généralisant 2.8. sous les conditions dites, et en déduire le

\uline{Corollaire} 2.\deleted{9}\handwritten{10}. $\underline C$, $\underline C'$ étant comme dans 2.8. , la catégorie des
foncteurs exacts à \deleted{droite} gauche $T$ de $\underline C$ dans $\underline C'$ est équivalente à la
catégorie des $A$-$\underline C'$-modules $H$ qui sont ``localement annulés par des
puissances de $I$'', i.e. qui sont limites inductives des sous-modules
$H_n=\operatorname{Hom}_A(A_n,H)$. On obtient deux foncteurs quasi-inverses par $T\handwritten{\leadsto}\underset{\handwritten{\longrightarrow}}{\lim}_{\handwritten{n}}T(A_n\unsure{)}$
et par $H\handwritten{\leadsto}T_H$, où $T_H(M)=\operatorname{Hom}_A(M,H)$ .

\uline{Remarque} 2.1\deleted{0}\handwritten{1}. Lorsqu'on ne suppose rien sur l'existence de limites
inductives dans $\underline C'$, on obtient encore une équivalence entre $\underline{\operatorname{Hom}}(\underline C^o,\underline C')_g$
et la catégorie des systèmes inductifs ``co-$\deleted{\unsure{J}}\handwritten{I}$-adiques'' $(H_n)$ de $A$-\inserted{$C'$-}modules
i.e. des systèmes inductifs de tels modules tels que \inserted{\handwritten{a)}} chaque $H_n$ soit
annulé par $I^{n+1}$ (donc en fait un $A_n$-$\underline C'$-module) et b) pour $m\handwritten{\geq}n$, l'homo%
% archive sheet number: 9
